\chapter{KẾT LUẬN VÀ HƯỚNG PHÁT TRIỂN}
\label{ch:chuong5}

\section{Kết luận}
\label{sec:ket_luan}

Sau thời gian nghiên cứu, thiết kế và tiến hành thực nghiệm nghiêm túc, đề tài \textbf{``Hệ thống Quản lý Ký túc xá Thông minh tích hợp IoT và Học máy (SmartDorm)''} đã hoàn thành xuất sắc toàn bộ các mục tiêu đặt ra ban đầu, mang lại những kết quả khoa học và ứng dụng thực tiễn nổi bật sau:

\begin{enumerate}
    \item \textbf{Về mặt phần cứng IoT và thiết bị nhúng:}
    \begin{itemize}
        \item Đã thiết kế và lắp ráp thành công mạch điều khiển nhúng trên vi điều khiển ESP32 DevKit V1 hoạt động bền bỉ, tích hợp đầy đủ cảm biến vi khí hậu DHT22, cảm biến rò rỉ khí gas/khói MQ-2, đầu đọc thẻ từ RFID RC522 13.56MHz và module relay 4 kênh cách ly quang.
        \item Giải quyết triệt để bài toán an toàn điện học đường bằng phương pháp ước tính công suất thông minh số hóa (Digital Proxy) thay vì sử dụng cảm biến điện áp lưới 220V trực tiếp, loại bỏ hoàn toàn nguy cơ chập cháy và điện giật trong môi trường ký túc xá.
    \end{itemize}

    \item \textbf{Về mặt kiến trúc phần mềm và truyền thông:}
    \begin{itemize}
        \item Xây dựng thành công kiến trúc 5 tầng chuẩn mực, kết hợp giao thức truyền thông nhẹ MQTT (QoS 1) với độ trễ truyền gói tin dưới $40\text{ ms}$ và WebSockets (Socket.io) cập nhật dữ liệu thời gian thực xuống trình duyệt dưới $15\text{ ms}$.
        \item Thiết kế cơ sở dữ liệu quan hệ tối ưu hóa chuỗi thời gian PostgreSQL tích hợp TimescaleDB với 12 bảng thực thể Prisma ORM chặt chẽ.
        \item Áp dụng các mẫu thiết kế kinh điển: Observer Pattern cho tiếp nhận sự kiện cảm biến, Strategy Pattern trong phân loại cảnh báo, Command Pattern với khả năng Hoàn tác lệnh (Undo Last Command), và Rule Engine tự động ngắt điện an toàn PCCC khi rò rỉ khí gas $\ge 80\%$.
    \end{itemize}

    \item \textbf{Về mặt Trí tuệ nhân tạo (Machine Learning):}
    \begin{itemize}
        \item Huấn luyện thành công mô hình Champion Gradient Boosting Machine với hàm tổn thất Huber Loss trên 13 đặc trưng, đạt hiệu năng dự báo phụ tải điện năng vượt trội ($R^2 = 0.9607$, $\text{MAE} = 8.28\text{ kW}$, $\text{MAPE} = 5.12\%$).
        \item Triển khai thuật toán Isolation Forest phát hiện bất thường thiết bị điện với chỉ số F1-Score đạt $94.6\%$ và độ trễ suy luận siêu thấp $11.8\text{ ms}$.
        \item Tích hợp trợ lý ảo thông minh AI Copilot hỗ trợ giải đáp trực tiếp cho sinh viên và ban quản lý.
    \end{itemize}

    \item \textbf{Về mặt ứng dụng và kiểm thử:}
    \begin{itemize}
        \item Phát triển ứng dụng web Next.js 16 giao diện Dark Mode hiện đại, trực quan, phân quyền chặt chẽ giữa Quản trị viên (Admin Portal) và Sinh viên (Student Portal).
        \item Đạt tỷ lệ $100\%$ thành công trên 84 ca kiểm thử tự động toàn diện từ Unit Test, Integration Test đến E2E Test, khẳng định tính ổn định cao và sẵn sàng đưa vào vận hành thử nghiệm.
    \end{itemize}
\end{enumerate}

\section{Những hạn chế của đề tài}
\label{sec:han_che}

Bên cạnh những kết quả tích cực đã đạt được, hệ thống vẫn tồn tại một số điểm hạn chế cần tiếp tục nghiên cứu hoàn thiện:
\begin{itemize}
    \item \textbf{Phạm vi phủ sóng không dây:} Hệ thống hiện tại sử dụng sóng Wi-Fi 2.4GHz để truyền thông từ ESP32 về máy chủ. Trong các tòa nhà KTX có kết cấu bê tông dày, vùng phủ sóng Wi-Fi có thể bị suy hao tại một số góc phòng khuất.
    \item \textbf{Quy mô thực nghiệm phần cứng:} Mô hình phần cứng hiện mới được chế tạo và thử nghiệm ở quy mô phòng mẫu (Prototype). Cần tiếp tục đóng gói công nghiệp hoàn thiện hộp bảo vệ (Enclosure) và kiểm định chuẩn chống cháy trước khi lắp đặt đại trà hàng loạt phòng.
    \item \textbf{Tính thích nghi của mô hình AI:} Mô hình Gradient Boosting hiện được huấn luyện trên tập dữ liệu tổng hợp chuẩn hóa. Khi triển khai tại một cơ sở ký túc xá mới với các thói quen sinh hoạt và vùng khí hậu khác biệt, mô hình cần có thời gian thu thập dữ liệu mới để tái huấn luyện (Retraining).
\end{itemize}

\section{Hướng phát triển trong tương lai}
\label{sec:huong_phat_trien}

Nhằm nâng cao hơn nữa giá trị thực tiễn và tính hiện đại của hệ thống SmartDorm, nhóm tác giả đề xuất các hướng phát triển tiếp theo như sau:
\begin{enumerate}
    \item \textbf{Mở rộng giao thức truyền thông diện rộng (LoRaWAN / Zigbee):} Tích hợp chuẩn truyền thông LoRaWAN hoặc Zigbee Mesh vào các nút mạng cảm biến phòng, giúp tiết kiệm năng lượng, xuyên tường tốt hơn và giảm phụ thuộc vào hạ tầng Wi-Fi của tòa nhà.
    \item \textbf{Ứng dụng Thị giác máy tính (AI Computer Vision):} Tích hợp camera AI nhận diện khuôn mặt (Face Recognition) tại cửa ra vào KTX nhằm xác thực sinh viên không cần tiếp xúc, kết hợp kiểm tra thẻ từ RFID tạo thành cơ chế xác thực đa nhân tố (MFA) chống mượn thẻ gian lận.
    \item \textbf{Triển khai TinyML trên biên (Edge Computing):} Nén và nạp trực tiếp các mô hình phát hiện bất thường siêu nhẹ lên vi điều khiển ESP32 (sử dụng TensorFlow Lite for Microcontrollers), cho phép phát hiện sự cố chập điện và ngắt tải ngay tại chỗ mà không cần chờ kết nối mạng về máy chủ.
    \item \textbf{Phát triển ứng dụng di động di động nguyên bản (Mobile Native App):} Xây dựng ứng dụng di động trên nền tảng React Native / Flutter cho cả hai hệ điều hành iOS và Android, tích hợp thông báo đẩy (Push Notifications) và quét mã QR Code nhận diện thiết bị phòng.
\end{enumerate}
\newpage
